\section{WAL：记录格式、写入线程、分段与恢复}

对应实现：\fpath{newdb/src/wal_manager.cpp}、\fpath{newdb/src/wal_codec.cpp}

\subsection{关键代码段：append\_record 如何分配 LSN 并计算校验}

\begin{lstlisting}[language=C++,caption={append\_record：构造 header + payload + checksum（节选）}]
hdr.magic = 0x57414C30;
hdr.lsn = next_lsn_.fetch_add(1) + 1;
hdr.txn_id = txn_id;
hdr.payload_len = (uint32_t)payload.size();
hdr.type = (uint8_t)op;
hdr.checksum = compute_checksum(&hdr, payload.data(), (uint32_t)payload.size());
\end{lstlisting}

要点：

\begin{itemize}
  \item \textbf{LSN 分配是并发安全的}：\texttt{next\_lsn\_} 是原子计数器，多线程调用也能得到全序 LSN。
  \item \textbf{checksum 覆盖范围}：header（除 checksum 字段）+ payload。恢复时先读 header/payload，再验证 checksum，失败则停止回放并报错。
\end{itemize}

\subsection{记录组成：Header + Payload + Checksum}

每条 WAL 记录由固定头与可选 payload 组成：

\begin{itemize}
  \item \textbf{magic}：常量 \texttt{0x57414C30}，用于快速识别；
  \item \textbf{lsn}：单调递增的日志序列号；
  \item \textbf{txn\_id}：事务 id；
  \item \textbf{payload\_len}：payload 长度；
  \item \textbf{type}：操作类型（INSERT/UPDATE/DELETE/COMMIT/ROLLBACK/CHECKPOINT/...）；
  \item \textbf{checksum}：CRC16-CCITT，覆盖 header(除 checksum 字段) 与 payload。
\end{itemize}

其中 checksum 计算使用 CCITT 多项式，作为轻量一致性检查。

\subsection{Payload 编码：table + row fields}

payload 的二进制布局由 \texttt{wal\_codec} 定义：

\begin{itemize}
  \item \texttt{u16 table\_name\_len} + table name bytes
  \item 若为行操作（INSERT/UPDATE/DELETE）：
    \begin{itemize}
      \item \texttt{u32 row\_id}
      \item \texttt{u32 row\_payload\_len}
      \item row payload bytes（由 tuple codec 编码得到）
    \end{itemize}
\end{itemize}

\subsubsection{关键代码段：payload 编码（table + row id + row payload）}

\begin{lstlisting}[language=C++,caption={walcodec::build\_payload 的布局（节选）}]
append_u16_le((uint16_t)table.size(), out);
out.insert(out.end(), table.begin(), table.end());
if (row_id != nullptr) {
    append_u32_le(*row_id, out);
    append_u32_le((uint32_t)row_payload->size(), out);
    out.insert(out.end(), row_payload->begin(), row_payload->end());
}
\end{lstlisting}

这说明 WAL payload 是“轻量自描述”的：只编码 table name 与可选的 row 字段；row 的内部结构由 heap payload（tuple codec）负责。

\subsection{写入模型：队列 + writer thread}

\texttt{WalManager::open()} 会启动 writer 线程：

\begin{itemize}
  \item 前台线程调用 \texttt{append\_record}/\texttt{flush}/\texttt{commit}/\texttt{rollback}；
  \item 操作被封装为 \texttt{QueuedWalOp} 进入队列；
  \item writer 线程串行写入文件，并通过 promise/future 把结果返回给调用方。
\end{itemize}

这样做的好处是：写入顺序天然序列化，且调用方可以同步等待“写入完成/刷盘完成”的结果。

\subsubsection{关键代码段：enqueue\_and\_wait（同步等待写入结果）}

\begin{lstlisting}[language=C++,caption={入队并等待 writer 返回 Status（节选）}]
auto done = std::make_shared<std::promise<Status>>();
auto fut = done->get_future();
op.done = done;
{
    std::lock_guard<std::mutex> lk(queue_mu_);
    queue_.push_back(std::move(op));
}
queue_cv_.notify_one();
return fut.get();
\end{lstlisting}

该模式把 writer 线程内部错误（例如 \texttt{fwrite} 失败）可靠地传播到调用方，避免“写入失败但上层以为成功”。

\subsection{刷盘策略（durability）}

\texttt{flush\_durable} 支持不同 sync 模式：

\begin{itemize}
  \item \textbf{Off}：仅 \texttt{fflush}；
  \item \textbf{Normal}：按时间间隔节流 \texttt{fsync}/\texttt{\_commit}；
\end{itemize}

在 Windows 上使用 \texttt{\_commit(\_fileno(fp))}，在类 Unix 上使用 \texttt{fsync(fileno(fp))}。

\subsubsection{关键代码段：Normal 模式的节流刷盘}

\begin{lstlisting}[language=C++,caption={flush\_durable：Normal 模式的时间间隔判断（节选）}]
if (sync_mode_ == WalSyncMode::Normal) {
    uint64_t now_ms = monotonic_ms();
    if (last_durable_flush_ms_ != 0 &&
        now_ms - last_durable_flush_ms_ < normal_sync_interval_ms_) {
        return Status::Ok();
    }
    last_durable_flush_ms_ = now_ms;
}
\end{lstlisting}

这是一种“吞吐优先”策略：在高频写入下减少 \texttt{fsync/\_commit} 次数，但会扩大崩溃丢失窗口（取决于 interval）。

\subsection{分段（segment rotate）}

当 WAL 文件大小超过阈值时，自动滚动到新段：

\begin{itemize}
  \item 基础段为 \texttt{db.wal}
  \item 后续段为 \texttt{db.wal.000001}、\texttt{db.wal.000002} ...
\end{itemize}

阈值通过环境变量 \texttt{NEWDB\_WAL\_SEGMENT\_BYTES} 设置。

\subsection{恢复（recovery）}

恢复逻辑的关键点：

\begin{itemize}
  \item 可为每个段建立稀疏 offset 索引（每 128 条采样）以支持按最小 LSN seek；
  \item 按事务聚合：同一事务内同一 \texttt{row\_id} 的多次写覆盖保存最后一次；
  \item COMMIT 时统一 apply 到 \texttt{HeapTable}，ROLLBACK 则丢弃；
  \item 最终 \texttt{rebuild\_indexes(schema)} 以获得查询索引。
\end{itemize}

\subsubsection{关键代码段：按事务聚合，commit 时统一 apply}

\begin{lstlisting}[language=C++,caption={recover：PendingTxnOps + COMMIT apply（节选）}]
// 事务内对同一 row_id 的操作，保留最后一次
PendingTxnOps& pending = txn_ops[hdr.txn_id];
pending.by_row_id[row.id] = { op, std::move(row) };

// COMMIT：对该 txn 的所有最终行操作执行 apply
for (auto& [row_id, op_row] : it->second.by_row_id) {
    WalOp op = op_row.first;
    Row& row = op_row.second;
    if (op == INSERT || op == UPDATE) { upsert_into_table(row); }
    if (op == DELETE) { append_tombstone(row.id); }
}
txn_ops.erase(it);
\end{lstlisting}

该实现的含义是：\textbf{事务原子性通过“commit 时才生效”实现}；在 commit 之前读到 WAL 的中间状态也不会被 apply 到表中。

可用环境变量：

\begin{itemize}
  \item \texttt{NEWDB\_RECOVER\_ENABLE\_OFFSET\_SEEK=1}
  \item \texttt{NEWDB\_RECOVER\_MIN\_LSN=<u64>}
\end{itemize}

